\textbf{Despeje del punto neutro.} La pendiente directa es $C_{m\alpha}=a_w(h-h_{ac})-K_t(h_t-h)+C_{m\alpha,f}$. Agrupar términos en $h$ deja $(a_w+K_t)h-a_wh_{ac}-K_th_t+C_{m\alpha,f}$. Igualar a cero produce la expresión de $h_n$; conservar $K_t$ en el denominador es necesario porque la cola también sustenta.\begin{enumerate}

\item La pendiente efectiva de cola incluye área, presión dinámica relativa y estela.

\[\fitmath{\mathrm{Kt}=0.93\cdot 0.2\cdot 4.3\cdot \left(1-0.33\right)\approx 0.535866\,}\]

\item Sumar las contribuciones a la sustentación total.

\[\fitmath{\mathrm{atot}=5+\mathrm{Kt}}\]

\[\fitmath{\mathrm{atot}=5+(0.535866)\approx 5.535866\,}\]

\item Imponer pendiente de momento nula y despejar la posición de CG.

\[\fitmath{\mathrm{hn}=\frac{5\cdot 0.25+\mathrm{Kt}\cdot 3.3-0.09}{\mathrm{atot}}}\]

\[\fitmath{\mathrm{hn}=\frac{5\cdot 0.25+(0.535866)\cdot 3.3-0.09}{(5.535866)}\approx 0.5289792\,}\]

\item Restar el margen mínimo exigido al punto neutro.

\[\fitmath{\mathrm{hmax}=\mathrm{hn}-0.08}\]

\[\fitmath{\mathrm{hmax}=(0.5289791697)-0.08\approx 0.4489792\,}\]

\item Comprobar la pendiente en el límite de CG.

\[\fitmath{\mathrm{Cm}=-\mathrm{atot}\cdot 0.08}\]

\[\fitmath{\mathrm{Cm}=-(5.535866)\cdot 0.08\approx -0.4428693\,\mathrm{1/rad}}\]

\end{enumerate}
